\documentclass[aspectratio=169,12pt]{beamer}
% Estilo canónico para las presentaciones de Métodos Numéricos 2026-2.
\usepackage[T1]{fontenc}
\usepackage[utf8]{inputenc}
\usepackage[spanish,es-nodecimaldot]{babel}
\usepackage{lmodern}
\usepackage{amsmath,amssymb,mathtools}
\usepackage{booktabs}
\usepackage{graphicx}
\usepackage{xcolor}
\usepackage{hyperref}

\definecolor{upazul}{HTML}{003B70}
\definecolor{upazulclaro}{HTML}{EAF2F8}
\definecolor{updorado}{HTML}{C5962E}
\definecolor{upgris}{HTML}{4A5560}
\definecolor{upverde}{HTML}{167C5A}
\definecolor{uprojo}{HTML}{B33A3A}

\setbeamercolor{normal text}{fg=black,bg=white}
\setbeamercolor{structure}{fg=upazul}
\setbeamercolor{frametitle}{fg=white,bg=upazul}
\setbeamercolor{title}{fg=upazul}
\setbeamercolor{subtitle}{fg=upgris}
\setbeamercolor{block title}{fg=white,bg=upazul}
\setbeamercolor{block body}{fg=black,bg=upazulclaro}
\setbeamercolor{alerted text}{fg=uprojo}
\setbeamercolor{example text}{fg=upverde}

\setbeamerfont{title}{size=\Huge,series=\bfseries}
\setbeamerfont{subtitle}{size=\Large}
\setbeamerfont{frametitle}{size=\LARGE,series=\bfseries}
\setbeamerfont{normal text}{size=\large}
\setbeamerfont{block title}{series=\bfseries}

\setbeamertemplate{navigation symbols}{}
\setbeamertemplate{itemize item}{\color{updorado}$\blacktriangleright$}
\setbeamertemplate{itemize subitem}{\color{upazul}$\bullet$}
\setbeamertemplate{blocks}[rounded][shadow=false]

\setbeamertemplate{frametitle}{%
  \nointerlineskip
  \begin{beamercolorbox}[wd=\paperwidth,ht=0.78cm,dp=0.24cm,leftskip=0.55cm]{frametitle}
    \usebeamerfont{frametitle}\insertframetitle
  \end{beamercolorbox}}

\setbeamertemplate{footline}{%
  \leavevmode
  \hbox{%
    \begin{beamercolorbox}[wd=.82\paperwidth,ht=2.7ex,dp=1.2ex,leftskip=.5cm]{author in head/foot}
      \color{upgris}\scriptsize Métodos Numéricos 2026-2 \hspace{0.8em} | \hspace{0.8em} Universidad del Pacífico
    \end{beamercolorbox}%
    \begin{beamercolorbox}[wd=.18\paperwidth,ht=2.7ex,dp=1.2ex,rightskip=.5cm plus1fil]{date in head/foot}
      \color{upazul}\scriptsize\hfill\insertframenumber/\inserttotalframenumber
    \end{beamercolorbox}}
  \vskip0pt}

\newcommand{\concepto}[1]{\textcolor{upazul}{\textbf{#1}}}
\newcommand{\resultado}[1]{\textcolor{upverde}{\textbf{#1}}}
\newcommand{\advertencia}[1]{\textcolor{uprojo}{\textbf{#1}}}
\newcommand{\codigo}[1]{\colorbox{upazulclaro}{\strut\texttt{#1}}}
\newcommand{\fuente}[1]{\vfill\hfill{\tiny\color{upgris}#1}}


\usepackage{pgfplots}
\pgfplotsset{compat=1.18}
\usetikzlibrary{babel}
\newcommand{\R}{\mathbb{R}}
\newcommand{\norm}[1]{\left\lVert#1\right\rVert}
\newcommand{\wegner}{\href{https://doi.org/10.1007/978-3-662-68697-3}{Wegner, \emph{Mathematische Einführung in Data Science}, cap. 17}}
\pgfplotsset{recta/.style={
  width=\linewidth,height=5.35cm, xmin=0,xmax=2.45, xlabel={$\alpha$},
  xtick={0,0.5,1,1.6,2.4}, xticklabels={$0$,$0.5$,$1$,$1.6$,$2.4$},
  axis lines=left, grid=major, grid style={upgris!15}, tick label style={font=\scriptsize},
  label style={font=\small}, legend style={font=\scriptsize,draw=none,fill=white},
  legend pos=north west, samples=121, clip mode=individual
}}
\title{Optimización sin restricciones}
\subtitle{Descenso del gradiente y búsqueda lineal}
\author{Oswaldo José Velásquez Castañón}
\institute{Universidad del Pacífico}
\date{Métodos Numéricos -- 2026-2}

\begin{document}

\begin{frame}[plain]
\begin{columns}[T,onlytextwidth]
\begin{column}{.73\textwidth}
\vspace{.55cm}{\usebeamerfont{title}\color{upazul}\inserttitle\par}
\vspace{.3cm}{\usebeamerfont{subtitle}\color{upgris}\insertsubtitle\par}
\vspace{.8cm}{\large\insertauthor\par}{\normalsize\insertdate\par}
\end{column}
\begin{column}{.22\textwidth}
\vspace{.55cm}\includegraphics[width=\linewidth]{UPlogo.pdf}
\end{column}
\end{columns}
\end{frame}

\section{Dirección de descenso}

\begin{frame}{Dirección y longitud de paso}
Buscamos minimizar una función diferenciable $f\colon\R^n\to\R$.
\[
g_k=\nabla f(x_k), \qquad
x_{k+1}=x_k+\alpha_k p_k,\qquad \alpha_k>0.
\]
Si $g_k^{\mathsf T}p_k<0$, un desplazamiento suficientemente corto en
dirección $p_k$ reduce $f$.
\begin{block}{Descenso del gradiente}
Con $p_k=-g_k$, la derivada direccional es $-\norm{g_k}^2$.
La dirección normalizada es la de descenso más rápido en norma euclidiana.
\end{block}
Estudiaremos primero un paso constante $\alpha$ y luego cómo elegir
$\alpha_k$ en cada iteración.
\end{frame}

\begin{frame}{Convexidad y optimalidad}
La convexidad exige que para $0\le t\le 1$,
\[
f((1-t)x+ty)\le (1-t)f(x)+tf(y).
\]
Si $f$ es diferenciable, equivale a que todo plano tangente sea soporte inferior.
\[
 f(y)\ge f(x)+\nabla f(x)^{\mathsf T}(y-x).
\]
\begin{block}{Consecuencia}
En una función convexa diferenciable, $\nabla f(x^*)=0$ caracteriza un
mínimo global. En general, un mínimo local diferenciable solo exige
estacionariedad como condición necesaria.
\end{block}
\end{frame}

\begin{frame}{El tamaño de paso en una cuadrática}
Para $f(x)=x^2$, el descenso satisface $x_{k+1}=(1-2\alpha)x_k$.
Con $x_0\ne 0$ y paso constante,
\begin{center}
\begin{tabular}{ll}
\toprule
Paso & Comportamiento \\
\midrule
$0<\alpha<1$ & Converge a $0$ \\
$\alpha=1$ & Alterna entre $x_0$ y $-x_0$ \\
$\alpha>1$ & Diverge en magnitud \\
\bottomrule
\end{tabular}
\end{center}
\medskip
También puede fallar un paso variable que disminuye muy rápido.
Con $\alpha_k=1/[2(k+2)^2]$, el producto de los factores
$1-2\alpha_k$ tiende a $1/2$.
\fuente{\wegner}
\end{frame}

\section{Convergencia con paso fijo}

\begin{frame}{Suavidad y modelo cuadrático superior}
El gradiente es $L$-Lipschitz si
$\norm{\nabla f(y)-\nabla f(x)}\le L\norm{y-x}$.
Con $s=y-x$, integramos la derivada sobre el segmento.
\[
 f(x+s)-f(x)-\nabla f(x)^{\mathsf T}s
 =\int_0^1[\nabla f(x+ts)-\nabla f(x)]^{\mathsf T}s\,dt
 \le \tfrac L2\norm{s}^2 .
\]
\begin{block}{Descenso con paso fijo}
Al usar $s=-\alpha\nabla f(x)$,
\[
 f(x+s)\le f(x)-\alpha(1-L\alpha/2)\norm{\nabla f(x)}^2 .
\]
Para $0<\alpha<2/L$ hay descenso cuando $\nabla f(x)\ne 0$.
\end{block}
\end{frame}

\begin{frame}{Convexidad y suavidad combinadas}
Si $f$ es convexa y $L$-suave,
\[
 f(u)-f(v)-\nabla f(v)^{\mathsf T}(u-v)
 \ge \frac{1}{2L}\norm{\nabla f(u)-\nabla f(v)}^2.
\]
\begin{block}{Demostración}
Sea $h(z)=f(z)-\nabla f(v)^{\mathsf T}z$. Es convexa y $v$ la minimiza.
El modelo superior con paso $1/L$ da
\[
 h(v)\le h\bigl(u-\nabla h(u)/L\bigr)
       \le h(u)-\frac{1}{2L}\norm{\nabla h(u)}^2 .
\]
\end{block}
Al sustituir $h$ se obtiene la desigualdad.
\fuente{\wegner}
\end{frame}

\begin{frame}{Convergencia sublineal con paso \texorpdfstring{$1/L$}{1/L}}
Supongamos que $f$ es convexa y $L$-suave y que existe $x^*$.
Sea $\delta_k=f(x_k)-f(x^*)$ y $R=\norm{x_0-x^*}$.
\begin{enumerate}
\item Al expandir $\norm{x_k-x^*-g_k/L}^2$ y aplicar la desigualdad anterior
con $u=x^*$, la distancia a $x^*$ decrece en sentido no estricto.
\item Por convexidad, $\delta_k\le R\norm{g_k}$. Por suavidad,
$\delta_k-\delta_{k+1}\ge\norm{g_k}^2/(2L)$.
\item Mientras $\delta_{k+1}>0$, esto da
$1/\delta_{k+1}-1/\delta_k\ge 1/(2LR^2)$.
\end{enumerate}
\[
\delta_k\le \frac{2LR^2}{k},\qquad k\ge 1.
\]
Una subsucesión acotada converge a un minimizador. La distancia a ese
minimizador decrece, de modo que converge toda la sucesión.
\fuente{\wegner}
\end{frame}

\begin{frame}{Convexidad fuerte y suavidad}
$f$ es $\mu$-fuertemente convexa si $f(x)-\mu\norm{x}^2/2$ es convexa
para $\mu>0$. Junto con la suavidad se obtiene
\[
 \tfrac\mu2\norm{y-x}^2
 \le f(y)-f(x)-\nabla f(x)^{\mathsf T}(y-x)
 \le \tfrac L2\norm{y-x}^2 .
\]
\begin{itemize}
\item La cota inferior implica coercividad y, por continuidad, existe un mínimo.
La convexidad estricta garantiza que sea único.
\item Ambas cotas implican $\mu\le L$.
Si $f$ es $C^2$, equivalen a $\mu I\preceq\nabla^2f(x)\preceq LI$.
\end{itemize}
El cociente $\kappa=L/\mu$ compara la curvatura máxima con la mínima.
\fuente{\wegner}
\end{frame}

\begin{frame}{Desigualdad de interpolación de Wegner}
Con $s=y-x$ y $d=\nabla f(y)-\nabla f(x)$,
\[
 d^{\mathsf T}s\ge
 \frac{\mu L}{\mu+L}\norm{s}^2+\frac{1}{\mu+L}\norm{d}^2 .
\]
Para $L>\mu$, $h(x)=f(x)-\mu\norm{x}^2/2$ es convexa y su modelo
superior tiene constante $L-\mu$. Aplicamos la desigualdad anterior
a $h$ en ambos órdenes y sumamos.
\[
 (d-\mu s)^{\mathsf T}s\ge\frac{\norm{d-\mu s}^2}{L-\mu}.
\]
Al desarrollar el cuadrado se obtiene el resultado.
Si $L=\mu$, las cotas coinciden y $d=\mu s$.
\fuente{\wegner, lema 17.22.}
\end{frame}

\begin{frame}{El paso que cancela el término del gradiente}
Sea $e_k=x_k-x^*$ y $g_k=\nabla f(x_k)$. Como $\nabla f(x^*)=0$,
la desigualdad anterior aplicada a $x^*$ y $x_k$ da
\[
\begin{aligned}
 \norm{e_{k+1}}^2
 &= \norm{e_k}^2-2\alpha g_k^{\mathsf T}e_k+\alpha^2\norm{g_k}^2\\
 &\le \left(1-\frac{2\alpha\mu L}{\mu+L}\right)\norm{e_k}^2
      +\left(\alpha^2-\frac{2\alpha}{\mu+L}\right)\norm{g_k}^2.
\end{aligned}
\]
\begin{block}{Elección de Wegner}
Con $\alpha=2/(\mu+L)$ desaparece el último término y queda
\[
 \norm{e_{k+1}}\le q\norm{e_k},\qquad q=\frac{L-\mu}{L+\mu}.
\]
\end{block}
\fuente{\wegner}
\end{frame}

\begin{frame}{Tasa lineal y condicionamiento}
\[
 \norm{x_k-x^*}\le q^k\norm{x_0-x^*},\qquad
 f(x_k)-f(x^*)\le\tfrac L2 q^{2k}\norm{x_0-x^*}^2 .
\]
La segunda cota usa el modelo superior en $x^*$.
\begin{center}
\begin{tabular}{rcc}
\toprule
$\kappa=L/\mu$ & $q$ en distancia & $q^2$ en la cota del objetivo \\
\midrule
$2$ & $1/3$ & $1/9$ \\
$8$ & $7/9$ & $49/81$ \\
$100$ & $99/101$ & $(99/101)^2$ \\
\bottomrule
\end{tabular}
\end{center}
\medskip
Con $\kappa$ grande, la garantía de reducción por iteración es pequeña.
¿Cómo elegir el paso cuando $\mu$ y $L$ son difíciles de estimar?
\end{frame}

\section{Búsqueda lineal}

\begin{frame}{La función sobre la dirección de descenso}
En la iteración actual fijamos $x$ y $p$, con $\nabla f(x)^{\mathsf T}p<0$.
\[
 \phi(\alpha)=f(x+\alpha p), \qquad
 \phi'(\alpha)=\nabla f(x+\alpha p)^{\mathsf T}p.
\]
Una búsqueda lineal inexacta evalúa candidatos hasta encontrar un paso
que cumpla las condiciones de aceptación.
\begin{block}{Ejemplo para las dos condiciones}
Tomamos $\phi(\alpha)=(\alpha-1)^2$, con $\phi(0)=1$,
$\phi'(0)=-2$ y mínimo en $\alpha=1$.
\end{block}
\end{frame}

\begin{frame}{Armijo y disminución suficiente}
\begin{columns}[c,onlytextwidth]
\begin{column}{.50\textwidth}
Para $0<c_1<1$ se exige
\[
 \phi(\alpha)\le\phi(0)+c_1\alpha\phi'(0).
\]
Con $c_1=0.2$, la recta de aceptación es $1-0.4\alpha$.
\medskip

El gráfico sombrea los pasos positivos que cumplen Armijo.
Aquí, $0<\alpha\le 1.6$.
\end{column}
\begin{column}{.49\textwidth}
\centering
\begin{tikzpicture}
\begin{axis}[recta,ymin=-.12,ymax=2.35,ylabel={Valor},ytick={0,0.5,1,1.5,2}]
\fill[upverde!18] (axis cs:0,-.12) rectangle (axis cs:1.6,2.35);
\addplot[upazul,very thick,domain=0:2.42] {(x-1)^2};
\addlegendentry{$\phi(\alpha)$}
\addplot[uprojo,dashed,thick,domain=0:2.42] {1-0.4*x};
\addlegendentry{Cota de Armijo}
\addplot[only marks,mark=*,upverde,mark size=2.2pt] coordinates {(1.6,0.36)};
\end{axis}
\end{tikzpicture}
\end{column}
\end{columns}
\end{frame}

\begin{frame}{Retroceso de Armijo}
Se fija un paso de prueba $\bar\alpha>0$ y un factor $0<\rho<1$.
Se prueban $\bar\alpha,\rho\bar\alpha,\rho^2\bar\alpha,\ldots$ hasta
que se cumpla Armijo.
\[
 \phi(\alpha)-\phi(0)-c_1\alpha\phi'(0)
   =(1-c_1)\alpha\phi'(0)+o(\alpha).
\]
Como $\phi'(0)<0$, los pasos positivos suficientemente pequeños se aceptan.
\begin{exampleblock}{En el gráfico anterior}
Con $\bar\alpha=2.4$ y $\rho=1/2$, se rechaza $2.4$ y se acepta $1.2$.
También se aceptaría $0.25$ si fuera el primer candidato.
\end{exampleblock}
\end{frame}

\begin{frame}{Powell--Wolfe y condición de curvatura}
\begin{columns}[c,onlytextwidth]
\begin{column}{.50\textwidth}
Con $0<c_1<c_2<1$, se pide Armijo y además
\[
 \phi'(\alpha)\ge c_2\phi'(0).
\]
Para $c_2=0.5$, se requiere $2(\alpha-1)\ge-1$.
La pendiente en $0.25$ es $-1.5$ y falla la segunda condición.
\medskip

\concepto{Ambas se cumplen en $[0.5,1.6]$.}
\end{column}
\begin{column}{.49\textwidth}
\centering
\begin{tikzpicture}
\begin{axis}[recta,ymin=-2.15,ymax=3.25,ylabel={Pendiente},ytick={-2,-1,0,1,2,3}]
\fill[upverde!18] (axis cs:.5,-2.15) rectangle (axis cs:1.6,3.25);
\addplot[upazul,very thick,domain=0:2.42] {2*(x-1)};
\addlegendentry{$\phi'(\alpha)$}
\addplot[uprojo,dashed,thick,domain=0:2.42] {-1};
\addlegendentry{Umbral $c_2\phi'(0)$}
\addplot[only marks,mark=*,upverde,mark size=2.2pt] coordinates {(.5,-1)};
\end{axis}
\end{tikzpicture}
\end{column}
\end{columns}
\end{frame}

\begin{frame}{Expansión y bisección para Powell--Wolfe}
\begin{enumerate}
\item Si el candidato cumple las dos condiciones, se acepta.
\item Si cumple Armijo pero la pendiente sigue por debajo del umbral,
se guarda como extremo inferior. Mientras falte un extremo superior,
se duplica el paso.
\item Si falla Armijo, se guarda como extremo superior. Con ambos
extremos, el próximo candidato es el punto medio. El extremo inferior
puede empezar en $0$.
\end{enumerate}
\begin{exampleblock}{Con $\bar\alpha=2.4$ en el ejemplo}
$2.4$ falla Armijo. La bisección de $[0,2.4]$ prueba $1.2$, que cumple ambas.
\end{exampleblock}
La prueba de curvatura requiere evaluar el gradiente.
El laboratorio registra ese costo y limita el número de intentos.
\end{frame}

\section{Laboratorio}

\begin{frame}{Ajuste de una curva de demanda}
\begin{columns}[c,onlytextwidth]
\begin{column}{.53\textwidth}
Con precios $P_i$ y cantidades $Q_i$ simulados, estimamos
$\widehat Q(P)=a+b(P-10)$.
\[
 \min_{a,b\in\R}\frac1{10}
   \sum_{i=1}^{5}\bigl[a+b(P_i-10)-Q_i\bigr]^2 .
\]
En $P=10$ la demanda ajustada es $a^*=48$.
Una unidad adicional de precio se asocia con $3.6$ unidades menos de demanda.
\end{column}
\begin{column}{.46\textwidth}
\centering
\begin{tikzpicture}
\begin{axis}[width=\linewidth,height=5.35cm,xmin=5,xmax=15,ymin=27,ymax=69,
  xlabel={Precio},ylabel={Cantidad},xtick={6,8,10,12,14},grid=major,
  grid style={upgris!15},tick label style={font=\scriptsize},label style={font=\small}]
\addplot[upazul,very thick,domain=6:14] {48-3.6*(x-10)};
\addplot[only marks,mark=*,mark size=2.6pt,updorado] coordinates
  {(6,63) (8,53) (10,49) (12,43) (14,32)};
\end{axis}
\end{tikzpicture}
\end{column}
\end{columns}
\end{frame}

\begin{frame}{El paso fijo en el ajuste de demanda}
Los precios centrados tienen media $0$ y segundo momento $8$. Por ello,
\[
 H=\begin{pmatrix}1&0\\0&8\end{pmatrix},\qquad
 \mu=1,\quad L=8,\quad \alpha=2/9.
\]
Con el paso fijo, el error en $a$ se multiplica por $7/9$ y el error en
$b$ por $-7/9$. La pendiente estimada alterna alrededor de $b^*=-3.6$.
\begin{block}{Comparación en el cuaderno}
Desde $(a_0,b_0)=(44,-1)$, se compara la distancia al ajuste de
referencia y el costo de las tres reglas de paso.
\end{block}
\end{frame}

\begin{frame}{Rosenbrock y protocolo de comparación}
\[
 f(x,y)=(1-x)^2+100(y-x^2)^2.
\]
El mínimo global está en $(1,1)$. El valle curvo permite observar un
comportamiento distinto al de la cuadrática de demanda.
\begin{itemize}
\item En cada problema se mantiene el mismo punto inicial y criterio
de parada para los tres métodos.
\item El cuaderno registra $f(x_k)-f(x^*)$, $\norm{\nabla f(x_k)}$ y las
evaluaciones de función y gradiente.
\item Los contornos muestran cuánto avanza cada trayectoria dentro del
mismo presupuesto de iteraciones.
\end{itemize}
\end{frame}

\begin{frame}{Material de laboratorio y referencias}
\begin{itemize}
\item Cuaderno \texttt{Optimizacion\_Descenso\_Gradiente\_2026-2.ipynb}.
  Incluye las dos gráficas de búsqueda lineal, el ajuste de demanda y
  la comparación en Rosenbrock.
\item \wegner. Descenso con paso fijo para funciones convexas.
\item \href{https://doi.org/10.2140/pjm.1966.16.1}{Armijo (1966)},
  \emph{Minimization of functions having Lipschitz continuous first partial derivatives}.
\item \href{https://doi.org/10.1137/1011036}{Wolfe (1969)},
  \emph{Convergence Conditions for Ascent Methods}.
\end{itemize}
\end{frame}

\end{document}
